% !TeX root = ../main.tex
\section{Style Transfer, Effect Representations and Sound Matching}
\label{sec:matching}

\subsection{Reference-based style transfer}

Style transfer of audio effects processes an input so that it sounds like a
reference. DeepAFx-ST predicts equaliser and compressor parameters from an
input and a reference, trained self-supervised with an audio loss; numerical
gradients (SPSA) let it work with non-differentiable effects at a cost in
stability~\cite{steinmetz2022style}. Grant trained the same way through legacy
plugins~\cite{grant2023style}. Koo et al.\ learn an effects encoder by
contrastive learning and render the input stems in the reference's mixing
style with a black-box network; its embedding distances on Demucs-separated
stems stayed close to those on clean stems~\cite{koo2023music}. Diff-MST
predicts the parameters of a differentiable mixing console for any number of
tracks~\cite{vanka2024diffmst}.

ST-ITO removes the need for a trained controller: CMA-ES searches the
parameters of an arbitrary chain of plugins to maximise the cosine similarity
between the AFx-Rep embeddings of the processed input and the
reference~\cite{steinmetz2024stito}. AFx-Rep is trained without dry references
to recognise which of 63 plugins and 10 presets produced an output, and it
beats CLAP, MFCCs and earlier effect encoders on zero-shot style
classification (0.86 against 0.56 for CLAP). The search uses a population of
64 for at most 25 steps and takes about a minute. The authors state that the
system does not work well for challenging cases such as guitar tone matching.
Yu et al.\ add a Gaussian prior fitted on real presets, which turns the search
into a maximum-a-posteriori estimate and reduces parameter error by up to a
third~\cite{yu2025improving}. ITO-Master found that optimising a learned
conditioning embedding beats optimising raw DSP parameters for
mastering~\cite{koo2025itomaster}.

TAPE predicts editable parameters of its own differentiable chain
(equaliser, compressor, waveshaper, a small residual network, reverberation)
from the processed reference in one pass, and runs causally in real
time~\cite{long2026tape}. Its comparison is the most instructive result for
our problem. CMA-ES over ST-ITO's chain with the AFx-Rep objective lowered the
embedding distance (0.259 against 0.660 for TAPE) but raised the spectral
error to 2.91, above the 2.34 of the unprocessed input, while TAPE reached
1.12. An embedding that is good for retrieval is not necessarily a good
objective to optimise.

\subsection{Effect representations}

Several encoders describe effects independently of content: AFx-Rep, the
effects encoder of Koo et al., Fx-Encoder++, the tone encoder of Chen et
al.~\cite{chen2024zeroshot}, the capture-trained encoder of Open-Amp
(91.1\% effect-class accuracy on EGDB~\cite{wright2025openamp}) and TAPE's.
Fx-Encoder++ extracts an instrument-wise effects embedding from a mixture
given an audio or text query~\cite{yeh2025fxencoder}. It also gives the only
measurement on separated stems: effect retrieval R@10 falls from 39.4\% on
true isolated stems to 20.5\% on Demucs-separated stems, which the authors
attribute to separation artefacts distorting effect cues. RelFx learns the
\emph{relative} effect between two processed clips, so it needs no dry
reference~\cite{liu2026beyond}; that matches the comparison between a stem and
the player's guitar through a candidate chain. General embeddings (OpenL3,
PANNs, CLAP) move monotonically with effect strength but not linearly, so they
are weak parameter estimators~\cite{deng2025investigating}. TimberAgent
retrieves the nearest of 1,063 guitar-effect presets with Gram matrices of
Wav2Vec2 features and reports the lowest parameter error among the retrieval
methods it tested, without claims on real audio~\cite{he2026timberagent}.

\subsection{Text and language models}

Text2FX optimises equaliser and reverberation parameters against a CLAP text
embedding~\cite{chu2025text2fx}; LLM2Fx predicts them directly with language
models, which for guitar equalisation did better than Text2FX
(distribution distance 0.15 against 0.42)~\cite{doh2025llm2fx}; FxSearcher uses
Bayesian optimisation over a plugin chain~\cite{ki2025fxsearcher}. The ``Does
it chug?'' dataset collects guitar-tone adjectives on clips that vary
equalisation and distortion~\cite{sutar2024chug}. None estimates amplifier or
drive settings, and text suits refinement (``more bite'') better than
estimation.

\subsection{Lessons from synthesizer sound matching}

Finding synthesizer parameters that reproduce a sound is the same inverse
problem with a better-studied history. Its findings carry over:
\begin{itemize}
\item \emph{Loss choice.} Training on audio losses alone often fails to
converge, especially for frequency-like parameters; pre-training on parameter
loss and then fine-tuning on spectral loss works best~\cite{masuda2021ddsp,uzrad2024diffmoog}.
LFO rates, delay times and wah positions are such parameters.
\item \emph{Optimisers.} On non-convex landscapes CMA-ES beat Adam (loss 2.09
against 4.60)~\cite{bogdan2026instrumental}; evolution strategies match
sounds through non-differentiable synthesizers with a CLAP
objective~\cite{cherep2024ctag}; reinforcement learning with rewards from
rendered audio improves out-of-domain matching~\cite{wu2026ddsynthrl}.
\item \emph{Ambiguity.} When different settings give the same sound, point
regression fails and generative models of the parameters do better~\cite{hayes2025equivariant}.
\item \emph{Proxies.} A neural proxy that maps presets to a pretrained audio
embedding makes black-box synthesizers trainable, at a resource
cost~\cite{combes2025proxies}.
\item \emph{Domain gap.} Every system degrades on real sounds outside its
training renders~\cite{hayes2024review}.
\end{itemize}
Taken together, the effect and synthesizer literatures agree on a hybrid:
predict an initial setting with a network trained on renders of the target
chain, refine it by gradient-free search through the real renderer, regularise
with a prior over realistic presets, and keep several candidates.
